\chapter{Loss of muscle mass}

\section*{Definition and Overview}
Loss of muscle mass, clinically categorised as muscle atrophy, sarcopenia, or cachexia depending on its underlying mechanism, refers to the progressive reduction in skeletal muscle bulk, strength, and function. While some decline in muscle mass is a natural part of the ageing process (primary sarcopenia), accelerated loss can occur secondary to physical inactivity, chronic illness, malignancy, or malnutrition. Muscle tissue plays a vital role not only in locomotion and physical stability but also in systemic metabolism, acting as a primary reservoir for amino acids and a major site for glucose disposal. Consequently, the loss of muscle mass is associated with profound clinical consequences, including metabolic dysfunction, physical frailty, an increased risk of falls and fractures, prolonged hospitalisations, and overall increased mortality.

\section*{Epidemiology}
The prevalence of muscle mass loss increases sharply with age. Primary sarcopenia affects approximately 5\% to 13\% of individuals aged 60 to 70 years, rising to as high as 50\% in those aged 80 years and older. Secondary muscle wasting, such as cachexia, is highly prevalent in specific patient populations, affecting up to 80\% of patients with advanced cancers (particularly upper gastrointestinal, pancreatic, and lung cancers), 15\% to 20\% of patients with chronic heart failure (cardiac cachexia), and a significant proportion of those with chronic obstructive pulmonary disease (COPD) and end-stage kidney disease. In Australia, the rising age of the population has made muscle wasting a major public health concern, contributing significantly to the burden of frailty and nursing home admissions.

\section*{Aetiology and Pathophysiology}
The loss of muscle mass results from an imbalance between protein synthesis (anabolism) and protein degradation (catabolism), driven by various physiological and pathological factors.
\begin{itemize}
    \item \textbf{Ageing-related cellular changes:} Progressive loss of alpha motor neurones in the spinal cord leads to denervation and subsequent atrophy of type II (fast-twitch) muscle fibres. Additionally, there is a decline in satellite cell (muscle stem cell) number and function.
    \item \textbf{Inflammatory cytokines:} In cachexia, chronic systemic inflammation driven by tumour- or immune-derived cytokines—such as tumour necrosis factor-alpha (TNF-$\alpha$), interleukin-6 (IL-6), and interferon-gamma—activates the ubiquitin-proteasome pathway, causing rapid skeletal muscle breakdown.
    \item \textbf{Disuse and immobilisation:} Prolonged bed rest, joint splinting, or sedentary lifestyle reduces mechanical loading on muscles, suppressing muscle protein synthesis and triggering rapid disuse atrophy.
    \item \textbf{Nutritional deficiencies:} Inadequate dietary protein or caloric intake, malabsorption syndromes, or anorexia lead to a negative nitrogen balance, forcing the body to catabolise skeletal muscle to maintain essential systemic amino acid pools.
    \item \textbf{Endocrine changes:} Declining levels of anabolic hormones, including testosterone, growth hormone, and insulin-like growth factor-1 (IGF-1), coupled with age-related insulin resistance and elevated cortisol, shift the metabolic balance toward muscle wasting.
    \item \textbf{Neurological disorders:} Conditions that impair motor innervation—such as amyotrophic lateral sclerosis (ALS), stroke, and peripheral neuropathies—prevent normal muscle activation, leading to neurogenic atrophy.
\end{itemize}

\section*{Clinical Features}
The clinical presentation of muscle mass loss depends on the rate of progression and the presence of underlying chronic diseases.
\begin{itemize}
    \item \textbf{Physical weakness:} Reduced grip strength and difficulty performing basic activities of daily living, such as carrying groceries, opening jars, or lifting objects.
    \item \textbf{Gait instability and slow walking speed:} An altered, slower, or unsteady gait, leading to difficulties crossing roads in a timely manner or climbing stairs.
    \item \textbf{Frequent falls:} Impaired balance and reduced lower-limb power increase the frequency of falls, often resulting in minor or major injuries.
    \item \textbf{Unintentional weight loss:} Visible thinning of the limbs, loosening of clothing, and a hollowed appearance of the temples and collarbones, especially prominent in cachectic states.
    \item \textbf{Profound fatigue:} Generalized lethargy, low energy levels, and exercise intolerance due to a reduction in active metabolic tissue and cardiopulmonary reserve.
    \item \textbf{Loss of functional independence:} Transition from independent living to requiring assistance with mobility, bathing, and dressing.
\end{itemize}

\section*{Differential Diagnosis}
Clinicians must investigate the underlying cause of muscle loss to distinguish between age-related changes, inflammatory wasting, and primary muscular or neurological diseases.
\begin{itemize}
    \item \textbf{Primary vs.\ secondary sarcopenia:} Differentiating age-related muscle decline from muscle loss driven by physical inactivity, poor diet, or specific organ failures.
    \item \textbf{Cachexia vs.\ simple starvation:} Starvation is characterised by fat loss with relative preservation of muscle mass, which can be reversed by feeding. Cachexia involves active inflammatory muscle degradation that cannot be fully reversed by nutritional supplementation alone.
    \item \textbf{Inflammatory myopathies:} Autoimmune disorders like polymyositis or dermatomyositis that present with proximal muscle weakness and elevated muscle enzymes.
    \item \textbf{Neuromuscular junction and motor neurone diseases:} Conditions such as myasthenia gravis, ALS, or severe peripheral neuropathies presenting with focal or generalised muscle wasting.
    \item \textbf{Endocrine muscle wasting:} Muscle weakness and wasting secondary to Cushing's syndrome (excess cortisol), severe hypothyroidism, or uncontrolled diabetes.
    \item \textbf{Drug-induced myopathy:} Muscle loss and weakness secondary to medications, most notably long-term corticosteroid therapy or statin-induced myotoxicity.
\end{itemize}

\section*{When to Seek Medical Care}
Patients should seek medical evaluation if they notice a progressive decline in physical strength, difficulty rising from a chair, recurrent falls, or significant, unintentional weight loss.

\textbf{Call 000 (or local emergency services) for:}
\begin{itemize}
    \item Acute, sudden-onset weakness in one or more limbs, which may indicate a stroke or acute spinal cord pathology.
    \item Respiratory distress or difficulty swallowing, suggesting severe weakness of the diaphragmatic or bulbar muscles.
    \item Inability to stand or move, or suspicion of a fracture (e.g. hip fracture) following a fall.
    \item Sudden collapse or altered state of consciousness.
\end{itemize}

\section*{Diagnosis and Investigations}
Evaluating muscle mass loss requires assessing muscle mass, muscle strength, and physical performance, alongside laboratory investigations to identify underlying causes.
\begin{itemize}
    \item \textbf{Dual-energy X-ray absorptiometry (DEXA):} The standard clinical tool used to measure appendicular lean mass and calculate the skeletal muscle mass index (SMI).
    \item \textbf{Handgrip dynamometry:} A simple, reliable bedside test used to measure maximum isometric grip strength; low grip strength ($\le 27$ kg for men, $\le 16$ kg for women) indicates probable sarcopenia.
    \item \textbf{Physical performance assessments:} Functional tests including gait speed measurement (threshold $\le 0.8$ m/s indicating impairment), the Timed Up and Go (TUG) test, and the Short Physical Performance Battery (SPPB).
    \item \textbf{Inflammatory and nutritional blood panel:} Full blood count, C-reactive protein (CRP), erythrocyte sedimentation rate (ESR), albumin, prealbumin, and vitamin D levels to assess nutritional status and systemic inflammation.
    \item \textbf{Hormonal evaluation:} Testing thyroid-stimulating hormone (TSH), morning cortisol, and free/total testosterone to identify endocrine causes of wasting.
    \item \textbf{Creatine kinase (CK) and renal markers:} Measured to rule out active muscle inflammation (myositis) and assess renal function (creatinine, eGFR).
\end{itemize}

\section*{Management}
Management requires a personalised, multidisciplinary approach focusing on physical exercise, nutritional optimization, and treating the primary medical driver.
\begin{itemize}
    \item \textbf{Progressive resistance training (PRT):} The cornerstone of therapy. High-intensity resistance exercise (e.g. weight lifting, resistance bands) performed 2 to 3 times per week is the most effective intervention to stimulate muscle protein synthesis and hypertrophy.
    \item \textbf{Nutritional optimization:} Ensuring adequate protein intake of 1.2 to 1.5 g/kg of body weight per day in older adults, rich in essential amino acids (particularly leucine), distributed evenly across meals.
    \item \textbf{Vitamin D supplementation:} Supplementation in patients with deficiency to support muscle contractility and bone health, reducing fall risk.
    \item \textbf{Anabolic agents and hormonal replacement:} Under strict medical supervision, testosterone replacement therapy in hypogonadal men, or growth hormone replacement in confirmed deficiency states.
    \item \textbf{Appetite stimulants:} In cancer- or chronic disease-related cachexia, short-term use of corticosteroids or progestins (e.g. megestrol acetate) to stimulate appetite.
    \item \textbf{Physical and occupational therapy:} Graded exercise programs, balance training, and home assessments to prescribe assistive devices and improve safety.
\end{itemize}

\section*{Complications}
The loss of muscle mass leads to a cascade of physical and metabolic complications that severely compromise quality of life.
\begin{itemize}
    \item \textbf{Falls and fractures:} Reduced muscle power impairs the ability to recover from a slip, leading to falls and subsequent fragility fractures (e.g. hip, wrist, or vertebral fractures).
    \item \textbf{Loss of mobility and bedridden state:} Progressive weakness can lead to a complete inability to walk, resulting in pressure injuries, joint contractures, and hypostatic pneumonia.
    \item \textbf{Metabolic syndrome and insulin resistance:} Since skeletal muscle is the largest site for insulin-stimulated glucose uptake, its loss promotes insulin resistance, glucose intolerance, and type 2 diabetes.
    \item \textbf{Poor surgical and chemotherapy outcomes:} Reduced physiological reserve (frailty) leads to poor wound healing, increased surgical complications, and higher toxicities during cancer treatments.
    \item \textbf{Respiratory failure:} Atrophy of the diaphragm and intercostal muscles leads to hypoventilation, poor cough reflex, and recurrent respiratory infections.
\end{itemize}

\section*{Prognosis and Outcomes}
The prognosis for individuals with muscle mass loss is highly dependent on the stage at diagnosis and the reversibility of the underlying cause. Age-related sarcopenia and disuse atrophy show excellent response rates, with significant recovery of muscle volume and strength achievable through structured resistance exercise and nutritional support within 3 to 6 months. In contrast, disease-related cachexia carries a poorer prognosis; muscle wasting in advanced cancer or end-stage organ failure is often refractory to standard interventions, serving as a strong independent predictor of shortened survival.

\section*{Prevention and Self-Care}
Preventing muscle loss is a lifelong process that involves maintaining physical activity and nutritional adequacy.
\begin{itemize}
    \item \textbf{Lifelong strength training:} Incorporating resistance exercises targeting major muscle groups at least twice a week.
    \item \textbf{Protein-rich diet:} Consuming high-quality protein sources, such as lean meats, poultry, fish, eggs, dairy, tofu, and legumes, at every meal.
    \item \textbf{Avoiding prolonged sedentary periods:} Breaking up sitting time, mobilizing as soon as possible during hospital stays, and engaging in daily walking.
    \item \textbf{Adequate hydration:} Drinking plenty of water to support cellular hydration and metabolic function in muscle tissues.
    \item \textbf{Regular health check-ups:} Monitoring body weight and muscle function, and discussing any unexplained changes with a general practitioner.
\end{itemize}

\section*{Sources and Further Reading}
The following organisations and guidelines provide further information on sarcopenia and muscle wasting:
\begin{itemize}
    \item Healthdirect Australia. \textit{Sarcopenia (muscle loss) in older age.} https://www.healthdirect.gov.au/
    \item Australian and New Zealand Society for Sarcopenia and Frailty Research (ANZSFR). \textit{Sarcopenia diagnosis and management guidelines.} https://www.anzsfr.org.au/
    \item Mayo Clinic. \textit{Sarcopenia: What you need to know about age-related muscle loss.} https://www.mayoclinic.org/
    \item NHS. \textit{Keeping active and eating well to prevent muscle weakness.} https://www.nhs.uk/
    \item European Working Group on Sarcopenia in Older People (EWGSOP). \textit{Sarcopenia: European consensus on definition and diagnosis.} https://www.sarcopenia-ewgsop.org/
\end{itemize}
